\documentclass[12pt]{article}
\usepackage{graphicx} %package to manage images
\graphicspath{ {./figures/} }
\usepackage{caption}
\usepackage[font=scriptsize]{subcaption}
\captionsetup[figure]{labelsep=none}
\captionsetup[table]{labelsep=none}
\usepackage{bbm}
\usepackage{amsmath}
\usepackage{import}
\usepackage{array}
\usepackage{booktabs}     % no extra vertical space between paragraphs
\usepackage{indentfirst}
\usepackage{afterpage}
\usepackage{floatrow}
\usepackage{pdflscape}
\usepackage{soul}
\usepackage{float}
\usepackage{adjustbox}
\usepackage{longtable}
\usepackage{caption}
\usepackage{setspace}
\usepackage{afterpage}
\usepackage[margin=1in]{geometry}
\usepackage[round]{natbib}
\usepackage{hyperref}
\usepackage{titlesec}
\usepackage{threeparttable}
\usepackage{setspace}
\usepackage{float}
\floatstyle{plaintop}
\restylefloat{table}
\usepackage{rotating}
\usepackage{authblk}
\usepackage{makecell}


% Optional styling for nicer author/affil typography
\renewcommand\Authfont{\normalsize\bfseries}
\renewcommand\Affilfont{\small\normalfont}
\setlength{\affilsep}{0.3em} % space between author and affiliation

\setlength{\parindent}{0pt}   % standard indent (~0.25 inch)
\setlength{\parskip}{1em}       % no extra vertical space between paragraphs

\titleformat{\section}
  {\normalfont\bfseries\fontsize{14}{16}\selectfont}  % Bold 14pt
  {}{0pt}{}  % No numbering

\titleformat{\subsection}
  {\normalfont\bfseries\fontsize{12}{14}\selectfont}  % Bold 12pt
  {}{0pt}{}  % No numbering

\titleformat{\subsubsection}
  {\normalfont\itshape\fontsize{12}{14}\selectfont}  % Italic 12pt
  {}{0pt}{}  % No numbering
  
\begin{document}
\title{Fertility, Household Wealth, and Maternal Nutrition among Indian Social Groups}
\date{\today}
\maketitle






\begin{table}[!htbp]
    \centering
    \begin{threeparttable}[t]
        % \caption{:  Descriptive statistics by social group and pregnancy status}
        \label{tab:sumstats}

        \scriptsize
        \setlength{\tabcolsep}{2pt}
        \renewcommand{\arraystretch}{1.2}

        {
\def\sym#1{\ifmmode^{#1}\else\(^{#1}\)\fi}
\begin{tabular}{l*{2}{c}}
\hline\hline
            &\multicolumn{1}{c}{\shortstack{Underweight among \\ prepregnant women}}&\multicolumn{1}{c}{\shortstack{Underweight among \\ prepregnant women's \\ husbands}}\\
\hline
\hline
Adivasi n=10630&       28.66         &       14.39         \\
Dalit n=9962&       23.91         &       14.22         \\
OBC n=17901 &       21.74         &       11.16         \\
Forward n=8043&       12.79         &       5.618         \\
Muslim n=5590&       16.63         &       8.851         \\
\hline\hline
\end{tabular}
}

    \end{threeparttable}
\end{table}



\begin{table}[!htbp]
    \centering
    \begin{threeparttable}[t]
        % \caption{:  Descriptive statistics by social group and pregnancy status}
        \label{tab:sumstats}

        \scriptsize
        \setlength{\tabcolsep}{2pt}
        \renewcommand{\arraystretch}{1.2}

        {
\def\sym#1{\ifmmode^{#1}\else\(^{#1}\)\fi}
\begin{tabular}{l*{1}{c}}
\hline\hline
            &\multicolumn{1}{c}{Underweight among men \\\\ whose wife gave birth \\\\ in the last year}\\
\hline
\hline
Adivasi n=1225&       12.62         \\
Dalit n=1014&       13.53         \\
OBC n=1549  &       10.60         \\
Forward n=548&       7.116         \\
Muslim n=681&       9.093         \\
\hline\hline
\end{tabular}
}

    \end{threeparttable}
\end{table}

Just an exercise, I'd like to think about what a public policy article would look like for this, and what we could say if we weren't constrained by academic journal writing.

\begin{itemize}
    \item lack of caste specific nationally representative outcomes, upcoming caste census 2027
    \item lack of pregnancy monitoring/ longitudinal surveys of maternal health
    \item compare our results in context of Jaccha Baccha survey
    \item political climate of maternity entitlements
    \item would maternity entitlements help generally? and what about for closing gaps?
    \item what might a decomposition of pp underweight over time look like, what factors might be responsible for the improvement?
    \item why are health outcomes so disproportionately bad for Adivasi?
    \item what explains the slight advantage Muslim over forward caste in prepregnancy underweight?
    \item what unobserved variables would we ideally like to decompose by?
    \item what might prepregnancy underweight in SSA be now? is the gap observed in Coffey (2015) widening, the same, or getting smaller?
    \item what is different about intrahousehold status of women during pregnancy across social groups?
    \item quality of gestational duration measurement in NFHS
    \item what a “maternal health census” might look like 
\end{itemize}



\begin{table}[!htbp]
    \centering
    \begin{threeparttable}[t]
        \caption{:  Underweight prevalence of prepregnant women and men who's wife is pregnant}

        \scriptsize
        \setlength{\tabcolsep}{2pt}
        \renewcommand{\arraystretch}{1.2}

        \begin{tabular}{l>{\centering\arraybackslash}p{3cm}>{\centering\arraybackslash}p{3.6cm}*{2}{>{\centering\arraybackslash}p{3cm}}}
\toprule
Social Group & \shortstack{Underweight \\ (prepregnant women)$^1$} & \shortstack{Underweight \\ (husbands \\ of pregnant women)$^2$} & \shortstack{Eligible but \\ not measured (\%)$^3$} & Not at home (\%)$^4$ \\\\
\midrule
Adivasi&0.27 [0.25, 0.28]&0.16 [0.10, 0.21]&0.08&0.13\\
Dalit&0.23 [0.22, 0.24]&0.14 [0.10, 0.18]&0.14&0.24\\
OBC&0.20 [0.20, 0.21]&0.12 [0.09, 0.15]&0.12&0.28\\
Forward&0.16 [0.14, 0.17]&0.09 [0.04, 0.13]&0.13&0.21\\
Muslim&0.14 [0.13, 0.15]&0.12 [0.07, 0.16]&0.21&0.24\\
All five social groups&0.20 [0.19, 0.20]&0.12 [0.10, 0.14]&0.14&0.24\\
\bottomrule
\end{tabular}


        \begin{tablenotes}[para,flushleft]
        \item $^1$ this is the main result, repeated
        \item $^2$ this is the prevalence of underweight among men who's wife is currently pregnant
        \item $^3$ this is the proportion of currently pregnant women who's husbands were eligible for BMI measurements, but weren't measured

        \item $^4$ this is the proportion of currently pregnant women in households eligible for male BMI measurement who's husband was not at home at the time of the survey (if he was home, he still may have been ineligible for BMI measurement, if he is not between the ages of 15-54 - husband's age data is missing for almost 85\% of these women, so eligibility was not imputed)
        
        \end{tablenotes}
    \end{threeparttable}
\end{table}


\begin{table}[!htbp]
    \centering
    \begin{threeparttable}[t]
        \caption{:  Underweight prevalence of prepregnant women and men who's wife is pregnant or $\le$ 1 year postpartum}

        \scriptsize
        \setlength{\tabcolsep}{2pt}
        \renewcommand{\arraystretch}{1.2}

        \begin{tabular}{l>{\centering\arraybackslash}p{3cm}>{\centering\arraybackslash}p{3.6cm}*{2}{>{\centering\arraybackslash}p{3cm}}}
\toprule
Social Group & \shortstack{Underweight \\ (prepregnant women)$^1$} & \shortstack{Underweight \\ (husbands of women \\ pregnant or \\ $\le$1 year postpartum)$^2$} & \shortstack{Eligible but \\ not measured (\%)$^3$} & Not at home (\%)$^4$ \\\\
\midrule
Adivasi&0.28 [0.27, 0.29]&0.13 [0.11, 0.16]&0.09&0.14\\
Dalit&0.25 [0.24, 0.26]&0.14 [0.11, 0.16]&0.13&0.24\\
OBC&0.23 [0.23, 0.24]&0.11 [0.09, 0.13]&0.13&0.27\\
Forward&0.18 [0.17, 0.19]&0.08 [0.05, 0.10]&0.14&0.23\\
Muslim&0.17 [0.16, 0.18]&0.10 [0.07, 0.12]&0.18&0.25\\
All five social groups&0.22 [0.22, 0.23]&0.11 [0.10, 0.12]&0.14&0.24\\
\bottomrule
\end{tabular}


        \begin{tablenotes}[para,flushleft]
        \item $^1$ this is the main result, repeated
        \item $^2$ this is the prevalence of underweight among men who's wife is currently pregnant or $\le$ 1 year postpartum
        \item $^3$ this is the proportion of women who are currently pregnant or $\le$ 1 year postpartum who's husbands were eligible for BMI measurements, but weren't measured

        \item $^4$ this is the proportion of women who are currently pregnant or $\le$ 1 year postpartum in households eligible for male BMI measurement who's husband was not at home at the time of the survey (if he was home, he still may have been ineligible for BMI measurement, if he is not between the ages of 15-54 - husband's age data is missing for almost 85\% of these women, so eligibility was not imputed)
        
        \end{tablenotes}
    \end{threeparttable}
\end{table}




\end{document}
